% ===================================================================
% 1. CORE DOCUMENT UTILITIES & PAGE LAYOUT
% ===================================================================
\usepackage[margin=0.5in]{geometry}   
\usepackage{graphicx}                 
\usepackage{float}                   
\usepackage{subcaption}
\usepackage{booktabs}                 
\usepackage{array}                    
\usepackage{multicol}                 
\usepackage{parskip}                  
\usepackage{titling}                  
%\setlength{\droptitle}{-1cm}          
\usepackage{titlesec}

% --- Blank Page Command ---
\newcommand*{\blankpage}{
  \begingroup
    \setlength{\topskip}{0pt}
    \setlength{\parskip}{0pt}
    \vspace*{\fill}
    \centering
    \nointerlineskip 
    This page was intentionally left blank. \par
    \vspace{\fill}
  \endgroup
}

% --- TODO ---
\newcommand{\todo}[1]{\textcolor{red}{\textbf{[TODO: #1]}}}


% ===================================================================
% 2. TYPOGRAPHY, ENCODING, TEXT & LISTS
% ===================================================================
\usepackage[utf8]{inputenc}
\usepackage{textcomp}
\usepackage{lipsum,lmodern}
\usepackage[svgnames]{xcolor}
\usepackage{soul}             
\usepackage[inline]{enumitem}         
\usepackage{gensymb}                  
\usepackage{xspace}                  

% --- Highlighted inline text ---
\newcommand{\hltext}[1]{\sethlcolor{gray!10}\hl{#1}}


% ===================================================================
% 3. MATHEMATICS, LOGIC, AND SYMBOLS
% ===================================================================
\usepackage{proof}             
\usepackage{amsmath, amssymb, mathtools} 
\usepackage{empheq}      
\usepackage{cancel}
\usepackage{mathrsfs}                   
\usepackage{stmaryrd}                    
\usepackage{mathpartir}                  
\usepackage{semantic}                    
\usepackage{proofpack}             

% --- Paired delimiters ---
\newcommand{\set}[1]{\left\{#1\right\}}
\newcommand{\abs}[1]{\left\lvert#1\right\rvert}
\newcommand{\norm}[1]{\left\lVert#1\right\rVert}
\newcommand{\paren}[1]{\left( #1 \right)}
\newcommand{\brak}[1]{[#1]}
\newcommand{\floor}[1]{\left\lfloor#1\right\rfloor}
\newcommand{\ceil}[1]{\left\lceil#1\right\rceil}

% --- Logic and Semantics ---
\newcommand{\Iof}[1]{\mathcal{I}(#1)}
\newcommand{\Ivof}[1]{\mathcal{I}_{\varphi}(#1)}
\newcommand{\sat}{\vDash}
\newcommand{\unsat}{\nvDash}
\newcommand{\satphi}{\vDash^{\varphi}}
\newcommand{\unsatphi}{\nvDash^{\varphi}}
\newcommand{\der}{\vdash}
\newcommand{\derC}{\vdash_{\mathcal{C}}}
\newcommand{\notder}{\nvdash}

% --- Math operators ---
\newcommand{\argmax}[1]{\mathop{\mathrm{arg\,max}}_{#1}\,}
\newcommand{\argmin}[1]{\mathop{\mathrm{arg\,min}}_{#1}\,}
\newcommand{\probd}[1]{\textbf{P}(#1)}

% --- Set and tuple utilities ---
\newcommand{\tuple}[1]{\left\langle #1 \right\rangle}
\newcommand{\restr}[2]{#1\!\upharpoonright_{#2}} 
\newcommand{\sint}[1]{\llbracket #1 \rrbracket}  
\newcommand{\inverse}[1]{\ensuremath{{#1}^{-1}}} 
\newcommand{\powerset}[1]{\ensuremath{\mathcal{P}(#1)}} 
\newcommand{\setbuilder}[2]{\ensuremath{\{\,#1 \mid #2\,\}}}

% --- LaTeX3 function syntax macros ---
\ExplSyntaxOn
\NewDocumentCommand{\union}{m}{\ensuremath{\clist_use:nn {#1} { \mathbin{\cup} }}}
\NewDocumentCommand{\intersection}{m}{\ensuremath{\clist_use:nn {#1} { \mathbin{\cap} }}}
\NewDocumentCommand{\compose}{m}{\ensuremath{\clist_use:nn {#1} { \mathbin{\circ} }}}
\NewDocumentCommand{\cartprod}{m}{\ensuremath{\clist_use:nn {#1} { \mathbin{\times} }}}
\ExplSyntaxOff

% --- Total and Partial Functions ---
\newcommand{\func}[3]{\ensuremath{#1 : #2 \to #3}}
\newcommand{\pfunc}[3]{\ensuremath{#1 : #2 \rightharpoonup #3}}

% --- Indexed set operations ---
\newcommand{\bigunion}[1]{\ensuremath{\bigcup_{i \in I} {#1}_{i}}}
\newcommand{\bigintersection}[1]{\ensuremath{\bigcap_{i \in I} {#1}_{i}}}

% --- Common math helpers ---
\newcommand{\identity}[1]{\ensuremath{\operatorname{Id}_{#1}}}
\newcommand{\bigo}[1]{\mathcal{O}(#1)}
\newcommand{\diff}{\mathop{}\!\mathrm{d}}


% ===================================================================
% 4. TIKZ AND DIAGRAMS
% ===================================================================
\usepackage{tikz}
\usetikzlibrary{arrows.meta, positioning, trees, shapes, matrix, calc, decorations.pathreplacing}
\tikzset{
  vertex/.style={circle, draw, minimum size=18pt, inner sep=0pt},
  edge/.style={-},
  arc/.style={->, >=Stealth}
}


% ===================================================================
% 5. CODE, LISTINGS, AND ALGORITHMS
% ===================================================================
\usepackage{algorithm, algorithmicx, algpseudocode} % Algorithm support
\usepackage{tcolorbox}
\tcbuselibrary{listingsutf8, skins, fitting, theorems, raster, magazine, breakable, xparse}

% --- Code block environment ---
\newtcblisting{codeblock}{
  colback=gray!10,
  colframe=black!40,
  listing only,
  boxrule=0.1pt,
  arc=2pt,
  left=0pt, right=0pt, top=0pt, bottom=0pt,
  listing options={
    breaklines=true,
    basicstyle=\ttfamily,
    literate={\$}{{\textdollar}}1
  }
}

% --- Code block with title ---
\newtcblisting{codeblocktitle}[1][]{
  colback=gray!10,
  colframe=black!40,
  colbacktitle=gray!30,
  coltitle=black,
  listing only,
  boxrule=0.1pt,
  arc=2pt,
  left=0pt, right=0pt, top=0pt, bottom=0pt,
  #1,
  listing options={
    breaklines=true,
    basicstyle=\ttfamily,
    literate={\$}{{\textdollar}}1
  }
}

% --- Captioned code block environment ---
\lstnewenvironment{codeblocklst}[1][]
{
  \lstset{
    backgroundcolor=\color{gray!10},
    frame=single,
    rulecolor=\color{black!10},
    framerule=0.1pt,
    breaklines=true,
    basicstyle=\ttfamily,
    literate={\$}{{\textdollar}}1,
    captionpos=b,
    #1
  }
}{}

% --- Inline code snippet ---
\newcommand{\inlinecode}[1]{%
  \tcbox[on line, fontupper=\ttfamily, size=small, boxrule=0.1pt, colframe=gray!50!black, colback=gray!10!white]{\detokenize{#1}}%
}


% ===================================================================
% 6. BIBLIOGRAPHY, HYPERLINKS & CLEVEREF SETUP
% ===================================================================
\usepackage[
backend=biber,
style=alphabetic,
sorting=ynt
]{biblatex}
\addbibresource{refs.bib}

\usepackage{hyperref}
\hypersetup{
  colorlinks,
  citecolor=teal,
  linkcolor=blue,
  urlcolor=magenta
}
\usepackage[nameinlink,capitalize]{cleveref}


% ===================================================================
% 7. CUSTOM STRUCTURES AND ANCHORS
% ===================================================================
% --- Section titles ---
\newcommand{\labeledsection}[2]{\section{#1}\label{#2}}
\newcommand{\minititle}[1]{\subsubsection*{#1}}

% --- Anchors and references ---
\newcommand{\refterm}[2]{\phantomsection \label{#2}{\textbf{#1}}}
\newcommand{\linkterm}[2]{\hyperref[#2]{#1}}

% --- Definition emphasizer ---
\newcommand{\defemph}[1]{\textbf{#1}}


% ===================================================================
% 8. DEFINITIONS, THEOREMS, AND MATHEMATICAL ENVIRONMENTS
% ===================================================================

% --- Section-scoped definitions ---
\newcounter{definition}[section]
\renewcommand{\thedefinition}{\thesection.\arabic{definition}}

\NewDocumentCommand{\definition}{o +m m}{%
  \refstepcounter{definition}%
  \phantomsection
  \par\noindent\textbf{Def~\thedefinition.} \IfValueT{#1}{\textbf{#1: }}#2%
  \if\relax\detokenize{#3}\relax\else\label{#3}\fi
  \par
}
\NewDocumentCommand{\Definition}{o +m m}{\definition[#1]{#2}{#3}}

% --- Subsection-scoped definitions ---
\newcounter{subdefinition}[subsection]
\renewcommand{\thesubdefinition}{\thesubsection.\arabic{subdefinition}}

\NewDocumentCommand{\subdefinition}{o +m m}{%
  \refstepcounter{subdefinition}%
  \phantomsection
  \par\noindent\textbf{Def~\thesubdefinition.} \IfValueT{#1}{\textbf{#1: }}#2%
  \if\relax\detokenize{#3}\relax\else\label{#3}\fi
  \par
}
\NewDocumentCommand{\Subdefinition}{o +m m}{\subdefinition[#1]{#2}{#3}}

% --- Shared theorem counter per section ---
\newcounter{statement}[section]
\renewcommand{\thestatement}{\thesection.\arabic{statement}}

% --- Linked theorem-type counters ---
\newcounter{theorem}
\renewcommand{\thetheorem}{\thestatement}
\newcounter{lemma}
\renewcommand{\thelemma}{\thestatement}
\newcounter{proposition}
\renewcommand{\theproposition}{\thestatement}
\newcounter{corollary}
\renewcommand{\thecorollary}{\thestatement}
\newcounter{axiom}
\renewcommand{\theaxiom}{\thestatement}
\newcounter{conjecture}
\renewcommand{\theconjecture}{\thestatement}

% --- Base style for theorem-like boxes ---
% --- Unified Formal Statement Style ---
% Consistent border and fill; the bold title conveys the semantic status.
\tcbset{
  statementbox/.style={
    enhanced,
    boxrule=0.7pt,
    arc=2mm,
    titlerule=0pt,
    fonttitle=\bfseries,
    coltitle=black,
    colframe=black!45!white,
    colbacktitle=black!12!white,
    colback=black!2.5!white,
    left=8pt,
    right=8pt,
    top=6pt,
    bottom=6pt,
    toptitle=4pt,
    bottomtitle=4pt,
    before skip=8pt,
    after skip=8pt,
    breakable
  }
}

\NewDocumentCommand{\Theorem}{o +m m}{%
  \stepcounter{statement}\refstepcounter{theorem}%
  \phantomsection
  \if\relax\detokenize{#3}\relax\else\label{#3}\fi
  \begin{tcolorbox}[
    statementbox,
    title={Theorem~\thetheorem\IfValueT{#1}{: #1}}
  ]
  #2
  \end{tcolorbox}%
}

\NewDocumentCommand{\Proposition}{o +m m}{%
  \stepcounter{statement}\refstepcounter{proposition}%
  \phantomsection
  \if\relax\detokenize{#3}\relax\else\label{#3}\fi
  \begin{tcolorbox}[
    statementbox,
    title={Proposition~\theproposition\IfValueT{#1}{: #1}}
  ]
  #2
  \end{tcolorbox}%
}

\NewDocumentCommand{\Lemma}{o +m m}{%
  \stepcounter{statement}\refstepcounter{lemma}%
  \phantomsection
  \if\relax\detokenize{#3}\relax\else\label{#3}\fi
  \begin{tcolorbox}[
    statementbox,
    title={Lemma~\thelemma\IfValueT{#1}{: #1}}
  ]
  #2
  \end{tcolorbox}%
}

\NewDocumentCommand{\Corollary}{o +m m}{%
  \stepcounter{statement}\refstepcounter{corollary}%
  \phantomsection
  \if\relax\detokenize{#3}\relax\else\label{#3}\fi
  \begin{tcolorbox}[
    statementbox,
    title={Corollary~\thecorollary\IfValueT{#1}{: #1}}
  ]
  #2
  \end{tcolorbox}%
}

\NewDocumentCommand{\Axiom}{o +m m}{%
  \stepcounter{statement}\refstepcounter{axiom}%
  \phantomsection
  \if\relax\detokenize{#3}\relax\else\label{#3}\fi
  \begin{tcolorbox}[
    statementbox,
    title={Axiom~\theaxiom\IfValueT{#1}{: #1}}
  ]
  #2
  \end{tcolorbox}%
}

\NewDocumentCommand{\Conjecture}{o +m m}{%
  \stepcounter{statement}\refstepcounter{conjecture}%
  \phantomsection
  \if\relax\detokenize{#3}\relax\else\label{#3}\fi
  \begin{tcolorbox}[
    statementbox,
    title={Conjecture~\theconjecture\IfValueT{#1}{: #1}}
  ]
  #2
  \end{tcolorbox}%
}

% --- Solution environment ---
\NewDocumentEnvironment{solution}{}{%
  \par\noindent\textit{Solution: }\par
}{%
  \par
}

% --- Exercise Macro ---
\definecolor{exercisecolor}{HTML}{2563EB}

\newcounter{exercise}[section]
\renewcommand{\theexercise}{\thesection.\arabic{exercise}}

\tcbset{
  exercisebox/.style={
    enhanced,
    breakable,
    frame hidden,
    boxrule=0pt,
    colback=white,
    borderline west={3pt}{0pt}{exercisecolor},
    left=12pt,
    right=0pt,
    top=0pt,
    bottom=0pt,
    before skip=12pt,
    after skip=12pt,
    coltitle=exercisecolor,
    fonttitle=\bfseries
  }
}

\NewDocumentEnvironment{exercise}{o m}{%
  \refstepcounter{exercise}%
  \phantomsection
  \if\relax\detokenize{#2}\relax\else\label{#2}\fi
  \begin{tcolorbox}[
    exercisebox,
    title={Exercise~\theexercise\IfValueT{#1}{: #1}}
  ]
}{%
  \end{tcolorbox}%
}

% --- Example Macro ---
\definecolor{examplecolor}{HTML}{EA580C}

\newcounter{example}[section]
\renewcommand{\theexample}{\thesection.\arabic{example}}

\tcbset{
  examplebox/.style={
    enhanced,
    breakable,
    frame hidden,
    boxrule=0pt,
    colback=white,
    borderline west={3pt}{0pt}{examplecolor},
    left=12pt,
    right=0pt,
    top=0pt,
    bottom=0pt,
    before skip=12pt,
    after skip=12pt,
    coltitle=examplecolor,
    fonttitle=\bfseries
  }
}

\NewDocumentEnvironment{example}{o m}{%
  \refstepcounter{example}%
  \phantomsection
  \if\relax\detokenize{#2}\relax\else\label{#2}\fi
  \begin{tcolorbox}[
    examplebox,
    title={Example~\theexample\IfValueT{#1}{: #1}}
  ]
}{%
  \end{tcolorbox}%
}

% --- Excursion Macro ---
\definecolor{excursioncolor}{HTML}{7C3AED}

\tcbset{
  excursionbox/.style={
    enhanced,
    breakable,
    frame hidden,
    boxrule=0pt,
    colback=white,
    borderline west={4pt}{0pt}{excursioncolor},
    left=13pt,
    right=0pt,
    top=0pt,
    bottom=0pt,
    before skip=16pt,
    after skip=16pt,
    coltitle=excursioncolor,
    fonttitle=\large\bfseries
  }
}

\NewDocumentEnvironment{excursion}{o m}{%
  \phantomsection
  \if\relax\detokenize{#2}\relax\else\label{#2}\fi
  \begin{tcolorbox}[
    excursionbox,
    title={Excursion\IfValueT{#1}{: #1}}
  ]
}{%
  \end{tcolorbox}%
}

% --- Proof Environment ---
\definecolor{proofcolor}{HTML}{475569}

\newtcolorbox{proofbox}{
    enhanced,
    frame hidden,
    boxrule=0pt,
    colback=white,
    borderline west={3pt}{0pt}{proofcolor},
    left=12pt,
    right=0pt,
    top=0pt,
    bottom=0pt,
    before skip=8pt,
    after skip=12pt,
    title=Proof,
    coltitle=proofcolor,
    fonttitle=\itshape
}

\newenvironment{Proof}[1][Proof]
{%
    \begin{proofbox}
}
{%
    \unskip\nobreak\hfill\penalty50\quad\ensuremath{\square}%
    \end{proofbox}%
}


% ===================================================================
% 9. CUSTOM BOXES
% ===================================================================

% --- Note boxes ---
\newcommand{\commandnote}[1]{%
\begin{tcolorbox}[
    standard jigsaw,
    breakable,
    colframe=gray!50!black,
    title=Note,
    fonttitle=\bfseries
]
#1
\end{tcolorbox}%
}

% --- Question ---
\newcommand{\questionbox}[1]{%
\begin{tcolorbox}[
    enhanced,
    attach boxed title to top left={xshift=5mm, yshift=-3mm, yshifttext=-1mm},
    colback=blue!4!white,
    colframe=blue!65!black,
    colbacktitle=blue!75!black,
    coltitle=white,
    title=Question,
    fonttitle=\bfseries,
    boxed title style={size=small, colframe=blue!85!black, arc=1.5mm},
    arc=2mm,
    boxrule=0.8pt
]
#1
\end{tcolorbox}%
}

% --- Answer ---
\newcommand{\answerbox}[1]{%
\begin{tcolorbox}[
    enhanced,
    breakable,
    attach boxed title to top left={xshift=5mm, yshift=-3mm, yshifttext=-1mm},
    colback=cyan!5!white,
    colframe=cyan!55!black,
    colbacktitle=cyan!70!black,
    coltitle=white,
    title=Answer,
    fonttitle=\bfseries,
    boxed title style={size=small, colframe=cyan!85!black, arc=1.5mm},
    arc=2mm,
    boxrule=0.8pt
]
#1
\end{tcolorbox}%
}

% --- Titled Box ---
\newcommand{\titledbox}[2]{%
\begin{tcolorbox}[
    colback=blue!3!white,
    colframe=blue!70!white,
    title=#1,
    fonttitle=\bfseries,
    arc=2mm
]
#2
\end{tcolorbox}%
}

% --- Red Box ---
\newcommand{\redbox}[1]{%
\begin{tcolorbox}[
    colback=red!5!white,
    colframe=red!75!black,
    arc=2mm,
    boxrule=0.8pt
]
#1
\end{tcolorbox}%
}

% --- Centered Titled Box ---
\newcommand{\centeredtitledbox}[2]{%
\begin{tcolorbox}[
    enhanced,
    attach boxed title to top center={yshift=-3mm, yshifttext=-1mm},
    colback=blue!3!white,
    colframe=blue!75!black,
    colbacktitle=white,
    coltitle=blue!80!black,
    title=#1,
    fonttitle=\bfseries,
    boxed title style={size=small, colframe=blue!75!black, boxrule=1pt, arc=1.5mm},
    arc=2mm,
    boxrule=0.8pt
]
#2
\end{tcolorbox}%
}

% --- Attention ---
\newcommand{\attentionbox}[1]{%
\begin{tcolorbox}[
    colback=red!5!white,
    colframe=red!75!black,
    coltitle=white,
    title=\textbf{Attention},
    fonttitle=\bfseries,
    arc=2mm,
    boxrule=0.8pt
]
#1
\end{tcolorbox}%
}

% --- Observation ---
\newcommand{\observationbox}[1]{%
\begin{tcolorbox}[
    enhanced,
    attach boxed title to top left={xshift=5mm, yshift=-3mm, yshifttext=-1mm},
    colback=violet!4!white,
    colframe=violet!60!black,
    colbacktitle=violet!75!black,
    coltitle=white,
    title=Observation,
    fonttitle=\bfseries,
    boxed title style={size=small, colframe=violet!85!black, arc=1.5mm},
    arc=2mm,
    boxrule=0.8pt
]
#1
\end{tcolorbox}%
}

% --- Idea ---
\newcommand{\ideabox}[1]{%
\begin{tcolorbox}[
    enhanced,
    attach boxed title to top left={xshift=5mm, yshift=-3mm, yshifttext=-1mm},
    colback=orange!6!white,
    colframe=orange!70!black,
    colbacktitle=orange!85!black,
    coltitle=white,
    title=Idea,
    fonttitle=\bfseries,
    boxed title style={size=small, colframe=orange!95!black, arc=1.5mm},
    arc=2mm,
    boxrule=0.8pt
]
#1
\end{tcolorbox}%
}

% --- Notation ---
\newcommand{\notationbox}[1]{%
\begin{tcolorbox}[
    enhanced,
    attach boxed title to top left={xshift=5mm, yshift=-3mm, yshifttext=-1mm},
    colback=gray!6!white,
    colframe=gray!60!black,
    colbacktitle=gray!75!black,
    coltitle=white,
    title=Notation,
    fonttitle=\bfseries,
    boxed title style={size=small, colframe=gray!85!black, arc=1.5mm},
    arc=2mm,
    boxrule=0.8pt
]
#1
\end{tcolorbox}%
}

% --- Problem ---
\newcommand{\problembox}[1]{%
\begin{tcolorbox}[
    enhanced,
    breakable,
    attach boxed title to top left={xshift=5mm, yshift=-3mm, yshifttext=-1mm},
    colback=red!35!black!4!white,
    colframe=red!45!black,
    colbacktitle=red!45!black,
    coltitle=white,
    title=Problem,
    fonttitle=\bfseries,
    boxed title style={size=small, colframe=red!60!black, arc=1.5mm},
    arc=2mm,
    boxrule=0.8pt
]
#1
\end{tcolorbox}%
}

% --- Solution ---
\newcommand{\solutionbox}[1]{%
\begin{tcolorbox}[
    enhanced,
    breakable,
    attach boxed title to top left={xshift=5mm, yshift=-3mm, yshifttext=-1mm},
    colback=green!5!white,
    colframe=green!45!black,
    colbacktitle=green!60!black,
    coltitle=white,
    title=Solution,
    fonttitle=\bfseries,
    boxed title style={size=small, colframe=green!80!black, arc=1.5mm},
    arc=2mm,
    boxrule=0.8pt
]
#1
\end{tcolorbox}%
}


% ===================================================================
% 10. CLEVEREF CONFIGURATION
% ===================================================================

% --- Cleveref names ---
\crefname{section}{section}{sections}
\Crefname{section}{Section}{Sections}
\crefname{subsection}{section}{sections}
\Crefname{subsection}{Section}{Sections}
\crefname{figure}{figure}{figures}
\Crefname{figure}{Figure}{Figures}
\crefname{table}{table}{tables}
\Crefname{table}{Table}{Tables}
\crefname{equation}{equation}{equations}
\Crefname{equation}{Equation}{Equations}
\crefname{algorithm}{algorithm}{algorithms}
\Crefname{algorithm}{Algorithm}{Algorithms}
\crefname{theorem}{theorem}{theorems}
\Crefname{theorem}{Theorem}{Theorems}
\crefname{lemma}{lemma}{lemmas}
\Crefname{lemma}{Lemma}{Lemmas}
\crefname{corollary}{corollary}{corollaries}
\Crefname{corollary}{Corollary}{Corollaries}
\crefname{proposition}{proposition}{propositions}
\Crefname{proposition}{Proposition}{Propositions}
\crefname{codeblock}{listing}{listings}
\Crefname{codeblock}{Listing}{Listings}
\crefname{appendix}{appendix}{appendices}
\Crefname{appendix}{Appendix}{Appendices}
\crefname{definition}{Def}{Defs}
\Crefname{definition}{Def}{Defs}
\crefname{subdefinition}{Def}{Defs}
\Crefname{subdefinition}{Def}{Defs}
\crefname{exercise}{exercise}{exercises}
\Crefname{exercise}{Exercise}{Exercises}
\crefname{example}{example}{examples}
\Crefname{example}{Example}{Examples}
\crefname{axiom}{axiom}{axioms}
\Crefname{axiom}{Axiom}{Axioms}
\crefname{conjecture}{conjecture}{conjectures}
\Crefname{conjecture}{Conjecture}{Conjectures}

% --- Reference shorthands ---
\newcommand{\secref}[1]{\cref{#1}}   \newcommand{\Secref}[1]{\Cref{#1}}
\newcommand{\figref}[1]{\cref{#1}}   \newcommand{\Figref}[1]{\Cref{#1}}
\newcommand{\tabref}[1]{\cref{#1}}   \newcommand{\Tabref}[1]{\Cref{#1}}
\newcommand{\eqrefc}[1]{\cref{#1}}   \newcommand{\Eqref}[1]{\Cref{#1}}
\newcommand{\defref}[1]{\cref{#1}}  \newcommand{\Defref}[1]{\Cref{#1}} 
\newcommand{\subdefref}[1]{\cref{#1}} \newcommand{\Subdefref}[1]{\Cref{#1}}
\newcommand{\thmref}[1]{\cref{#1}}   \newcommand{\Thmref}[1]{\Cref{#1}}
\newcommand{\lemref}[1]{\cref{#1}}   \newcommand{\Lemref}[1]{\Cref{#1}}
\newcommand{\corref}[1]{\cref{#1}}   \newcommand{\Corref}[1]{\Cref{#1}}
\newcommand{\propref}[1]{\cref{#1}}   \newcommand{\Propref}[1]{\Cref{#1}}
\newcommand{\axref}[1]{\cref{#1}}   \newcommand{\Axref}[1]{\Cref{#1}}
\newcommand{\conjref}[1]{\cref{#1}}  \newcommand{\Conjref}[1]{\Cref{#1}}
\newcommand{\exampleref}[1]{\cref{#1}} \newcommand{\Exampleref}[1]{\Cref{#1}}
\newcommand{\exerciseref}[1]{\cref{#1}} \newcommand{\Exerciseref}[1]{\Cref{#1}}
\newcommand{\lstref}[1]{\cref{#1}}   \newcommand{\Lstref}[1]{\Cref{#1}}  
\newcommand{\appref}[1]{\cref{#1}}   \newcommand{\Appref}[1]{\Cref{#1}}
\newcommand{\Algref}[1]{\Cref{#1}}
